% ------------------------------------------------------------------------
% file `main-exercise-8-exercise-body.tex'
%   in folder `xsim/'
%
%     exercise of type `exercise' with id `8'
%
% generated by the `exercise' environment of the
%   `xsim' package v0.21 (2022/02/12)
% from source `main' on 2022/06/13 on line 6
% ------------------------------------------------------------------------
    Voici, pour la production de l'année 2021, le relevé
    des longueurs des gousses de vanille d'un cultivateur
    de Tahaa.
    \begin{center}
        \begin{tblr}{hlines, vlines, columns={r, 1cm}, column{1}={l, 2.5cm}}
            Longueur (cm) & 12        & 15        & 17          & 22          & 23        \\
            Effectif      & \num{600} & \num{800} & \num{1 800} & \num{1 200} & \num{600}
        \end{tblr}
    \end{center}

    \begin{enumerate}
        \item Quel est l'effectif total de cette production ?
        \item Détermine l'étendue de cette série.
        \item Le cultivateur peut seulement conditionner les
              gousses de vanille dans des tubes de \SI{20}{cm} de
              long. Quel pourcentage de cette production a-t-il
              pu conditionner sans plier les gousses ?
        \item La chambre d'agriculture décerne une récompense
              (un \og label de qualité \fg{}) aux agriculteurs si :
              \begin{itemize}
                  \item  la longueur moyenne des gousses de leur
                        production est supérieure ou égale à \SI{16,5}{cm} ;
                  \item plus de la moitié des gousses de leur
                        production a une taille supérieure à \SI{17,5}{cm}.
              \end{itemize}
              Ce cultivateur pourra-t-il recevoir ce \og label de
              qualité \fg{} ?
    \end{enumerate}
